\documentclass[pdf, hyperref={unicode}, aspectratio=169]{beamer}

\usepackage{./diploma-slides-latex-template/styles}
\usepackage{booktabs}
\usepackage{listings}

\title{Разработка и исследование пулов потоков с механизмом work-stealing на языке C++}

\subtitle{Выпускная квалификационная работа студента СпецВО-магистратуры}

\pdfstringdefDisableCommands{
  \def\\{}
  \def\,{}
  \def\textbf#1{<#1>}
}

\author[Спиридонов Кирилл Анатольевич]
{
  \textbf{Студент группы М8О-214СВ-24:} Спиридонов Кирилл Анатольевич\\
  \ \textbf{Научный руководитель:} д.ф.-м.н., доцент Морозов Александр Юрьевич
  % Обратите внимание на пробел в начале строки
}

\institute[Московский авиационный институт]
{
  Московский авиационный институт (национальный исследовательский университет)\\
  Институт № 8 «Компьютерные науки и прикладная математика»\\
  Кафедра № 806 «Вычислительная математика и программирование» 
}

\date{Москва --- \the\year}

\logo{\includegraphics[height=1cm]{./diploma-slides-latex-template/img/mai}}

\begin{document}

\frame{\titlepage}

% ─── 1. Актуальность ───────────────────────────────────────────────────────────
\begin{frame}
\frametitle{Актуальность темы}

\begin{columns}[T]

\column{0.65\textwidth}
\begin{itemize}
  \item В начале 2000-х годов рост тактовой частоты процессоров достиг
        физических пределов~--- индустрия перешла к многоядерным
        архитектурам
  \item Флагманские серверные процессоры содержат \textbf{сотни ядер};
        параллелизм стал повседневной инженерной реальностью
  \item Классический пул потоков с \textbf{единой глобальной очередью}
        плохо масштабируется: очередь превращается в узкое место при
        большом числе потоков
  \item Стандартная библиотека \textbf{C++17 не предоставляет} готовой
        реализации пула потоков
\end{itemize}

\column{0.42\textwidth}
\begin{block}{Проблема}
  Отсутствует реализация пула потоков с механизмом work-stealing,
  построенная исключительно на стандарте C++17 без внешних
  зависимостей и подкреплённая систематическим экспериментальным
  исследованием
\end{block}

\end{columns}
\end{frame}


% ─── 2. Цель и задачи ──────────────────────────────────────────────────────────
\begin{frame}
\frametitle{Цель и задачи работы}

\textbf{Цель~---} разработать библиотеку на языке C++, реализующую
пул потоков с механизмом work-stealing, и экспериментально исследовать
её производительность в сравнении с классическим пулом потоков

\medskip
\textbf{Задачи:}
\begin{itemize}
  \item Проанализировать существующие подходы и реализации пулов потоков
  \item Спроектировать архитектуру библиотеки с единым интерфейсом
        для двух реализаций
  \item Реализовать \texttt{SimpleThreadPool} и
        \texttt{WorkStealingPool} без внешних зависимостей
  \item Верифицировать корректность посредством автоматических тестов
        и Thread Sanitizer
  \item Провести сравнительное исследование производительности
        и сформулировать практические рекомендации
\end{itemize}
\end{frame}


% ─── 3. Постановка задачи ──────────────────────────────────────────────────────
\begin{frame}
\frametitle{Постановка задачи}

\begin{columns}[T]

\column{0.5\textwidth}
\textbf{Дано:}
\begin{itemize}
  \item Многоядерная вычислительная система
  \item Поток задач с неизвестной заранее структурой нагрузки
  \item Теоретическая оценка Блюмофе--Лейзерсона:
        \[T_P \le \frac{T_1}{P} + O(T_\infty)\]
  \item Алгоритм Chase--Lev для двусторонней очереди
\end{itemize}

\column{0.55\textwidth}
\textbf{Требования к решению:}
\begin{itemize}
  \item Язык~--- C++17, без внешних зависимостей
  \item Единый интерфейс \texttt{submit()} для обеих реализаций
  \item Корректность при конкурентном доступе
        (отсутствие гонок данных)
\end{itemize}

\medskip
\textbf{Необходимо:}\\
Библиотека с двумя взаимозаменяемыми реализациями пула потоков
и экспериментальное сравнение их характеристик
\end{columns}
\end{frame}


% ─── 4. Теоретическая основа ───────────────────────────────────────────────────
\begin{frame}
\frametitle{Теоретическая основа: механизм work-stealing}

\begin{columns}[T]

\column{0.52\textwidth}
\textbf{Модель Блюмофе--Лейзерсона:}
\begin{itemize}
  \item Параллельное вычисление~--- ориентированный ациклический граф
  \item $T_1$~--- работа (однопроцессорное время)
  \item $T_\infty$~--- глубина (критический путь)
  \item Гарантия жадного work-stealing:
        \[T_P \le \frac{T_1}{P} + O(T_\infty)\]
\end{itemize}

\textbf{Закон Амдала:}
\[S_P \le \frac{1}{\alpha + \frac{1-\alpha}{P}}\]
Накладные расходы на синхронизацию увеличивают~$\alpha$

\column{0.45\textwidth}
\textbf{Принцип работы deque:}

\begin{block}{Владелец (LIFO, нижний конец)}
  \texttt{push(task)} \quad \texttt{pop(task)}\\
  В типичном случае~--- без конкуренции
\end{block}

\begin{block}{Вор (FIFO, верхний конец)}
  \texttt{steal(task)}\\
  Забирает «старые» крупные задачи
\end{block}

\medskip
Асимметрия обеспечивает низкую конкуренцию
и хорошую локальность данных

\end{columns}
\end{frame}


% ─── 5. Архитектура библиотеки ─────────────────────────────────────────────────
\begin{frame}
\frametitle{Архитектура библиотеки}

\begin{columns}[T]

\column{0.45\textwidth}
\textbf{Три независимых модуля:}

\medskip
\begin{block}{\texttt{WorkStealingDeque}}
  Двусторонняя очередь.\\
  Асимметричный доступ\\
  владельца и воров.\\
  Синхронизация: \texttt{std::mutex}
\end{block}

\begin{block}{\texttt{SimpleThreadPool}}
  Единая глобальная очередь.\\
  \texttt{std::mutex} +
  \texttt{std::condition\_variable}
\end{block}

\begin{block}{\texttt{WorkStealingPool}}
  По одной \texttt{WorkStealingDeque}\\
  на каждый поток.\\
  Случайный выбор жертвы (\texttt{mt19937})
\end{block}

\column{0.52\textwidth}
\textbf{Единый интерфейс:}

\medskip
\begin{semiverbatim}
template<typename F, typename... Args>
auto submit(F\&\& f, Args\&\&... args)
  -> std::future<
       invoke\_result\_t<F,Args...>>;
\end{semiverbatim}

\medskip
\textbf{Управление жизненным циклом:}
\begin{itemize}
  \item Конструктор: создаёт рабочие потоки
        (по умолчанию~---
        \texttt{hardware\_concurrency()})
  \item Деструктор: graceful shutdown~---
        дожидается завершения всех задач
  \item Классы некопируемы и неперемещаемы
\end{itemize}

\medskip
\textbf{Стек технологий:}
C++17, CMake~$\ge$~3.14,
Google Test~1.14.0,
Google Benchmark~1.8.3

\end{columns}
\end{frame}


% ─── 6. Ключевые алгоритмы ─────────────────────────────────────────────────────
\begin{frame}
\frametitle{Алгоритм рабочего цикла WorkStealingPool}

\begin{columns}[T]

\column{0.50\textwidth}
\textbf{Рабочий цикл потока с индексом $i$:}

\medskip
\begin{enumerate}
  \item \textbf{Своя очередь:} \texttt{queues[i].pop(task)}\\
        Если есть~--- выполнить, перейти к~1
  \item \textbf{Кража:} до 32 попыток\\
        случайно выбрать жертву $j \ne i$,\\
        \texttt{queues[j].steal(task)}\\
        Если удалось~--- выполнить, перейти к~1
  \item \textbf{Ожидание:}
        \texttt{cv.wait\_for(1\,мс)}\\
        Предикат: есть задачи \textbf{или} флаг
        остановки
  \item \textbf{Завершение:} если флаг
        остановки~--- опустошить свою очередь
        и выйти
\end{enumerate}

\column{0.47\textwidth}
\textbf{Нетривиальные проблемы реализации:}

\begin{alertblock}{Гонка в деструкторе SimpleThreadPool}
  \texttt{notify\_all()} до входа потока в
  \texttt{wait()} $\Rightarrow$ уведомление
  теряется.\\
  \textbf{Решение:} устанавливать флаг под
  мьютексом
\end{alertblock}

\begin{alertblock}{Переполнение \texttt{size\_t} в pop()}
  \texttt{bottom\_ - 1} при
  \texttt{bottom\_ == 0} $\Rightarrow$
  segfault.\\
  \textbf{Решение:} проверять
  \texttt{bottom == top} до вычитания
\end{alertblock}

Все дефекты обнаружены \textbf{Thread Sanitizer}

\end{columns}
\end{frame}


% ─── 7. Тестирование ───────────────────────────────────────────────────────────
\begin{frame}
\frametitle{Верификация корректности}

\begin{columns}[T]

\column{0.52\textwidth}
\textbf{28 автоматических тестов (Google Test):}

\medskip
\begin{tabular}{lc}
\toprule
Набор тестов & Тестов \\
\midrule
\texttt{WorkStealingDeque} & 12 \\
\texttt{SimpleThreadPool}  & 6  \\
\texttt{WorkStealingPool}  & 10 \\
\midrule
\textbf{Итого} & \textbf{28} \\
\bottomrule
\end{tabular}

\column{0.45\textwidth}
\textbf{Ключевые тесты:}
\begin{itemize}
  \item \texttt{SingleElementRaceCondition}~---
        10\,000 итераций: владелец и вор
        одновременно претендуют на
        последний элемент; задача
        выполняется ровно один раз
  \item \texttt{DestructorWaitsPending}~---
        деструктор дожидается всех задач
  \item \texttt{HeavyLoadAllTasks}~---
        100\,000 задач под нагрузкой
\end{itemize}
\end{columns}

\medskip
\begin{block}{Результат}
  Все 28 тестов пройдены.
  Thread Sanitizer не обнаружил ни одной гонки данных.
  Сборка с \texttt{-fsanitize=thread} интегрирована в CMake.
\end{block}
\end{frame}


% ─── 8. Условия эксперимента ───────────────────────────────────────────────────
\begin{frame}
\frametitle{Условия экспериментального исследования}

\begin{columns}[T]

\column{0.45\textwidth}
\textbf{Тестовая машина:}

\medskip
\begin{tabular}{ll}
\toprule
Параметр & Значение \\
\midrule
Логических ядер & 16 \\
Тактовая частота & 5102{,}71~МГц \\
Кэш L1 & 32~КиБ $\times$ 8 \\
Кэш L2 & 1024~КиБ $\times$ 8 \\
Кэш L3 & 16384~КиБ $\times$ 1 \\
ОС & Linux \\
Компилятор & Clang, C++17 \\
Оптимизация & \texttt{-O2} \\
\bottomrule
\end{tabular}

\column{0.52\textwidth}
\textbf{Два сценария нагрузки:}

\begin{block}{Сценарий 1: много мелких задач}
  10\,000 задач, каждая~--- атомарный
  инкремент ($\approx$5~нс).\\
  Измеряет накладные расходы планировщика.
\end{block}

\begin{block}{Сценарий 2: неравномерная нагрузка}
  10\,000 задач: половина~---
  атомарный инкремент,\\
  половина~--- \texttt{sleep\_for(200~мкс)}.\\
  Целевой сценарий для work-stealing.
\end{block}

\medskip
Число потоков: \textbf{1, 2, 4, 8, 16}.\\
Измерение: реальное время (\texttt{UseRealTime}).\\
Число итераций: автоматически (Google Benchmark).

\end{columns}
\end{frame}


% ─── 9. Результаты ─────────────────────────────────────────────────────────────
\begin{frame}
\frametitle{Результаты экспериментального исследования}

\begin{columns}[T]

\column{0.50\textwidth}
\textbf{Сценарий 1: много мелких задач}

\medskip
\begin{tabular}{ccc}
\toprule
Потоков & Simple, мс & WS, мс \\
\midrule
1  & 4{,}52 & 4{,}39 \\
2  & \textbf{3{,}76} & 4{,}86 \\
4  & 6{,}77 & 7{,}93 \\
8  & 14{,}2 & 14{,}8 \\
16 & 12{,}5 & 15{,}2 \\
\bottomrule
\end{tabular}

\medskip
\small Обе реализации \alert{деградируют} с ростом
числа потоков: накладные расходы планировщика
многократно превышают полезную работу

\column{0.47\textwidth}
\textbf{Сценарий 2: неравномерная нагрузка}

\medskip
\begin{tabular}{ccc}
\toprule
Потоков & Simple, мс & WS, мс \\
\midrule
1  & 1268 & 1266 \\
2  & 638  & 635  \\
4  & 317  & 317  \\
8  & \textbf{159}  & 237  \\
16 & \textbf{80{,}3} & 98{,}1 \\
\bottomrule
\end{tabular}

\medskip
\small При 16 потоках SimpleThreadPool
быстрее по реальному времени,
но WorkStealingPool потребляет
\alert{в 3{,}8 раза меньше CPU}
(2{,}75~мс против 10{,}4~мс)

\end{columns}
\end{frame}


% ─── 10. Анализ результатов ────────────────────────────────────────────────────
\begin{frame}
\frametitle{Анализ результатов: фундаментальный компромисс}

\begin{columns}[T]

\column{0.50\textwidth}
\textbf{Почему WorkStealingPool медленнее\\по реальному времени?}

\begin{enumerate}
  \item \texttt{WorkStealingDeque} использует
        \texttt{std::mutex}~--- операции
        владельца \textbf{не} lock-free.\\
        Теоретическое преимущество
        локальных очередей не достигается
  \item Механизм ожидания:
        \texttt{cv.wait\_for(1\,мс)}~---
        потоки засыпают при отсутствии
        задач; пробуждение с задержкой
        до 1~мс
  \item SimpleThreadPool пробуждает
        потоки \textbf{мгновенно} через
        \texttt{condition\_variable}
\end{enumerate}

\column{0.47\textwidth}
\textbf{Практические рекомендации:}

\begin{block}{SimpleThreadPool}
  Предпочтителен при жёстких
  требованиях к \textbf{латентности}:
  высоконагруженные серверы,
  интерактивные приложения
\end{block}

\begin{block}{WorkStealingPool}
  Предпочтителен при требованиях
  к \textbf{энергоэффективности}:
  фоновые вычисления, мобильные
  устройства, системы с ограниченным
  энергопотреблением.\\
  \medskip
  Таймаут \texttt{wait\_for} настраиваем:
  меньше~--- ниже латентность,
  больше~--- меньше CPU
\end{block}

\end{columns}
\end{frame}


% ─── 11. Выводы ────────────────────────────────────────────────────────────────
\begin{frame}
\frametitle{Выводы и результаты работы}

\begin{enumerate}
  \item \textbf{Спроектирована и реализована} модульная библиотека
        на чистом C++17 без внешних зависимостей:
        \texttt{SimpleThreadPool}, \texttt{WorkStealingPool},
        \texttt{WorkStealingDeque} с единым интерфейсом
        \texttt{submit()}

  \item \textbf{Верифицирована корректность:} 28 тестов пройдены,
        Thread Sanitizer не обнаружил гонок данных;
        выявлены и устранены нетривиальные дефекты

  \item \textbf{Установлен фундаментальный компромисс:}
        SimpleThreadPool~--- меньшая латентность (80{,}3~мс при
        16 потоках), WorkStealingPool~--- в 3{,}8 раза меньше
        CPU (2{,}75~мс против 10{,}4~мс) при неравномерной нагрузке

  \item \textbf{Научная новизна:} систематически исследовано влияние
        механизма ожидания на соотношение реального и процессорного
        времени; выявлены условия, при которых блокирующая
        синхронизация в deque нивелирует преимущества work-stealing
\end{enumerate}

\medskip
\begin{block}{Перспективы}
  Lock-free deque (алгоритм Chase--Lev + hazard pointers);
  NUMA-топологическая оптимизация;
  рекурсивные сценарии («разделяй и властвуй»)
\end{block}
\end{frame}


% ─── 12. Код ───────────────────────────────────────────────────────────────────
\begin{frame}
\frametitle{Исходный код библиотеки}

Полный исходный код размещён в открытом репозитории GitHub.

\begin{figure}
  \includegraphics[height=0.62\textheight]{./diploma-slides-latex-template/img/qr-code}
\end{figure}

\end{frame}

\end{document}